% -*- coding: utf-8 -*-
% !TEX program = xelatex
\documentclass[plain,noanswer,medmath]{../../template/jnuexam}
\usepackage{../../template/kysx}

\begin{document}


\examtitle{1993}{1}


\makepart{填空题}{本题共5小题，每小题3分，满分15分}


\begin{problem}
函数$F(x) = \int_1^x \left(2 - \frac{1}{\sqrt{t}} \right)\dt$（$x > 0$）的单调减少区间为 \fillin{}．
\end{problem}


\begin{problem}
由曲线$\begin{cases}
  3x^2 + 2y^2 = 12, \\
  z = 0
\end{cases}$绕$y$轴旋转一周得到的旋转面在点$(0,\sqrt{3},\sqrt{2})$处的指向外侧的单位法向量为 \fillin{}．
\end{problem}


\begin{problem}
设函数$f(x) = \pi x + x^2$（$- \pi < x < \pi$）的傅里叶级数展开式为
$$\frac{a_0}{2} + \sum_{n = 1}^{\infty}(a_n\cos nx + b_n\sin nx),$$
则其中系数$b_3$的值为 \fillin{}．
\end{problem}


\begin{problem}
设数量场$u = \ln \sqrt{x^2 + y^2 + z^2}$，则$\div(\grad u) =$ \fillin{}．
\end{problem}


\begin{problem}
设$n$阶矩阵$A$的各行元素之和均为零，且$A$的秩为$n - 1$，则线性方程组$Ax = 0$的通解为 \fillin{}．
\end{problem}


\makepart{选择题}{本题共5小题，每小题3分，满分15分}


\begin{problem}
设$f(x) = \int_0^{\sin x} \!\!\sin(t^2)\dt$，$g(x) = x^3 + x^4$，则当$x \to 0$时，$f(x)$是$g(x)$的\pickout{}
\begin{abcd}
\item 等价无穷小．				
\item 同阶但非等价无穷小． 
\item 高阶无穷小．				
\item 低阶无穷小．
\end{abcd}
\end{problem}


\begin{problem}
双纽线$(x^2 + y^2)^2 = x^2 - y^2$所围成的区域面积可用定积分表示为\pickout{}
\begin{abcd}
\item $2\int_0^{\frac{\pi}{4}} \cos 2\theta \d\theta$．					
\item $4\int_0^{\frac{\pi}{4}} \cos 2\theta \d\theta$．
\item $2\int_0^{\frac{\pi}{4}} \sqrt{\cos 2\theta} \d\theta$．				
\item $\frac{1}{2}\int_0^{\frac{\pi}{4}} (\cos 2\theta)^2 \d\theta$．
\end{abcd}
\end{problem}


\begin{problem}
设有直线$L_1:\frac{x - 1}{1} = \frac{y - 5}{-2} = \frac{z + 8}{1}$与$L_2:\begin{cases}
  x - y = 6 \\
  2y + z = 3
\end{cases}$，则$L_1$与$L_2$的夹角为\pickout{}
\begin{abcd}
\item $\frac{\pi}{6}$．        
\item $\frac{\pi}{4}$．        
\item $\frac{\pi}{3}$．        
\item $\frac{\pi}{2}$．
\end{abcd}
\end{problem}


\begin{problem}
设曲线积分$\int_L [f(x) - \e^x]\sin y\dx - f(x)\cos y\dy$与路径无关，
其中$f(x)$具有一阶连续导数，且$f(0) = 0$，则$f(x)$等于\pickout{}
\begin{abcd}
\item $\frac{\e^{- x} - \e^x}{2}$．   
\item $\frac{\e^x - \e^{-x}}{2}$．   
\item $\frac{\e^x + \e^{-x}}{2} - 1$．   
\item $1-\frac{\e^x + \e^{-x}}{2}$．
\end{abcd}
\end{problem}


\begin{problem}
已知$Q = \begin{pmatrix}
  1 & 2 & 3 \\
  2 & 4 & t \\
  3 & 6 & 9
\end{pmatrix}$，$P$为三阶非零矩阵，且满足$PQ = O$，则
\begin{abcd}
\item $t = 6$时，$P$的秩必为$1$．
\item $t = 6$时，$P$的秩必为$2$．
\item $t \ne 6$时，$P$的秩必为$1$．
\item $t \ne 6$时，$P$的秩必为$2$．
\end{abcd}
\end{problem}


\makepart{计算题}{本题共3小题，每小题5分，满分15分}


\begin{problem}
求 $\lim_{x \to \infty} \left(\sin \frac{2}{x} + \cos \frac{1}{x} \right)^x$．
\end{problem}


\begin{problem}
求 $\int \frac{x\e^x}{\sqrt{\e^x - 1}}\dx$．
\end{problem}


\begin{problem}
求微分方程$x^2 y' + xy = y^2$，满足初始条件$y|_{x = 1} = 1$的特解．
\end{problem}


\makepart{}{本题满分6分}

\begin{problem}
计算$\iint_{\Sigma}2xz\dy\dz + yz\dz\dx - z^2 \dx\dy$,
其中$\Sigma$是由曲面$z = \sqrt{x^2 + y^2}$与$z = \sqrt{2 - x^2 - y^2}$ 所围立体的表面外侧．
\end{problem}


\makepart{}{本题满分7分}

\begin{problem}
求级数$\sum_{n = 0}^{\infty}\frac{(- 1)^n(n^2 - n + 1)}{2^n}$的和．
\end{problem}


\makepart{证明题}{本题共2小题，每小题5分，满分10分}


\begin{problem}
设在$[0,+\infty)$上函数$f(x)$有连续导数，且$f'(x) \ge k > 0$, $f(0) < 0$，
证明$f(x)$在$(0,+\infty)$内有且仅有一个零点．
\end{problem}


\begin{problem}
设$b > a > \e$，证明$a^b > b^a$．
\end{problem}


\makepart{}{本题满分8分}

\begin{problem}
已知二次型$f(x_1,x_2,x_3) = 2x_1^2+ 3x_2^2+ 3x_3^2+ 2ax_2x_3$（$a > 0$）通过正交变换化成标准形$f = y_1^2+ 2y_2^2+ 5y_3^2$，求参数$a$及所用的正交变换矩阵．
\end{problem}


\makepart{}{本题满分6分}

\begin{problem}
设$A$是$n \times m$矩阵，$B$是$m \times n$矩阵，其中$n < m$,$E$是$n$阶单位矩阵，
若$AB = E$，证明$B$的列向量组线性无关．
\end{problem}


\makepart{}{本题满分6分}

\begin{problem}
设物体$A$从点$(0,1)$出发，以速度大小为常数$v$沿$y$轴正向运动．物体$B$从点$(- 1,0)$与$A$同时出发，其速度大小为$2v$，方向始终指向$A$，试建立物体$B$的运动轨迹所满足的微分方程，并写出初始条件．
\end{problem}


\makepart{填空题}{本题共2小题，每小题3分，满分6分}


\begin{problem}
一批产品共有$10$个正品和$2$个次品，任意抽取两次，每次抽一个，抽出后不再放回，
则第二次抽出的是次品的概率为 \fillin{}．
\end{problem}


\begin{problem}
设随机变量$X$服从$(0,2)$上的均匀分布，则随机变量$Y = X^2$在$(0,4)$内的概率分布密度$f_Y(y) =$ \fillin{}．
\end{problem}


\makepart{}{本题满分6分}

\begin{problem}
设随机变量$X$的概率分布密度为$f(x) = \frac{1}{2}\e^{- |x|}$, $-\infty < x < +\infty$．
\begin{enumerate}
\item 求$X$的数学期望$E(X)$和方差$D(X)$．
\item 求$X$与$|X|$的协方差，并问$X$与$|X|$是否不相关？
\item 问$X$与$|X|$是否相互独立？为什么？
\end{enumerate}
\end{problem}



\end{document}
